% Section 4.1 -- Conditional-rate reward (hero). Plain, direct, number-backed prose (RAG Results
% register). No clefts/padding, no editorial teasers.
\subsection{Conditional-rate reward}
\label{sec:reward}

The ideal reward would score each forecast against the true rate $\eta(x)$, but $\eta(x)$ is never
observed. The only label for a play is its realized outcome $y$, a single draw from that rate. A reward
built on the draw, $r=1-(p-y)^2$, is unbiased but noisy, since two identical favorites, one that wins
and one that loses, produce equal and opposite gradients. The noise does not average out with more plays
at a fixed state, because it is intrinsic to a binary label, and in the small groups that GRPO compares
\citep{grpo} it overwhelms the per-state signal.

We score each forecast against an estimate of the rate instead. We bin the training plays by score
margin, time remaining, and pregame point spread, and set each bin to the fraction of its plays the
possession team won. Bins with few plays are shrunk toward the overall rate by an empirical-Bayes prior,
so a sparse bin is not governed by its few plays \citep{efron_morris1975}. The estimate
$\hat p(x)$ uses only realized outcomes and the public pregame line, never the live market probability,
which we hold out for evaluation. The reward is
\begin{equation}
  r \;=\; 1 - \bigl(p - \hat p(x)\bigr)^2 .
  \label{eq:reward}
\end{equation}
The target uses no human labels. It comes from the same outcomes a binary reward would use, averaged
over states to yield a low-variance estimate of $\eta(x)$. On held-out 2024 games it scores a Brier of
$0.143$, close to the market's $0.136$, and it supplies a dense reward at every state. We call $\hat p$
the teacher, because a model trained on \cref{eq:reward} can be no better calibrated than $\hat p$
itself.

An ablation under identical training isolates the target as the cause (\cref{tab:reward}). Rewarding the
realized outcome reaches a Brier of $0.166$, but its calibration error drifts to $0.10$ as the policy
fits individual outcomes. An equal blend of outcome and rate is worse on both, $0.181$ and $0.121$,
because it reintroduces the outcome's noise. The rate alone is best on both, at $0.154$ and $0.050$.

\begin{table}[t]
\centering
% IN-TRAINING held-out (n=128), direct prediction, beta=0.01, scale_rewards on.
% NOT recomputable from test preds; transcribed from notes/experiments.md.
% p-hat & blend at lr 1e-5; outcome at lr 5e-6 (its most stable run).
\begin{tabular}{lcc}
\toprule
Reward target & Brier & ECE \\
\midrule
Realized outcome $y\in\{0,1\}$ & 0.166 & 0.10 \\
Blend $\tfrac{1}{2}y + \tfrac{1}{2}\hat{p}$ & 0.181 & 0.121 \\
Empirical rate $\hat{p}$ (ours) & 0.154 & 0.050 \\
\bottomrule
\end{tabular}

\caption{Reward-target ablation under identical training (direct prediction, \cref{sec:decoupling};
in-training held-out, $n=128$). The realized-outcome reward is higher-variance and its calibration
drifts; blending the outcome back in is worse on both axes; the conditional rate $\hat p$ is the most
accurate and the most stable.}
\label{tab:reward}
\end{table}
